\section{Methods}\label{sec:methods}

\subsection{Study design}
This is a retrospective evaluation of a general-purpose LLM with a Python sandbox, used as a zero-shot classifier of
12-lead ECGs on a public dataset; the model was blind to labels and PTB-XL identifiers. We use zero-shot to mean that
the model's parameters were not updated (no training or fine-tuning), that the prompt contained no labelled examples
and no class prevalences, and that no retrieval of external information was used. The prompt did define the superclasses as PTB-XL does, and
the configuration was developed on labelled validation records (``Development and freeze''). The input format,
prompt, model settings and analysis plan were fixed in a written protocol that was frozen before the test fold was
run. Analyses added after the test run are listed under ``Post hoc analyses'' below and labelled post hoc where they
are reported. Ethics is described at the end of this section.

\subsection{Data, labels and split}
We used PTB-XL version 1.0.3~\citep{Wagner2020,Wagner2022physionet} (\NRecordsTotal{} records from \NPatients{}
patients, 10\,s, 12 leads, recorded between October 1989 and June 1996 at 400\,Hz and resampled to 500\,Hz by the
dataset's authors; the dates in the metadata are shifted per patient). PTB-XL annotates each record with statements
of the SCP-ECG standard (Standard Communication Protocol for Computer-Assisted Electrocardiography) and gives each
diagnostic statement a likelihood from 0 to 100, inferred from keywords in the ECG report that express the
cardiologist's certainty (0 where the report gave no such information)~\citep{Wagner2020}. Labels and split follow the
public benchmark~\citep{Strodthoff2021}: every diagnostic SCP-ECG statement of a record (as marked in PTB-XL's
\texttt{scp\_statements} table) contributes its diagnostic superclass (NORM, MI, STTC, CD or HYP), irrespective of its
likelihood; records without any superclass are removed, which leaves \NWithSuperclass{} records. All test-fold records
are marked in PTB-XL as validated by a human annotator. A record can have several superclasses. Folds 1--8
(\NTrain{} records) form the benchmark training set, fold 9 (\NVal{} records) the validation set and fold 10
(\NTest{} records) the test set; no patient appears in more than one of these sets (Supplementary
Table~\ref{tab:characteristics}). All development runs used fold 9. Folds 1--8 were used only to fit the age-and-sex
control, to compute the class prevalences reserved for failed records, and for the labelled example files tried
during development. In this study, the test fold was used only for the frozen test run, the pre-specified evaluation
and, after the test run, the post hoc analyses.

\subsection{Input}
Each record was written to a CSV file with a time column in milliseconds and the 12 leads (I, II, III, aVR, aVL, aVF,
V1--V6) in millivolts with three decimals, one row per sample: the full 10-second recording at 500\,Hz (5{,}000 rows),
unfiltered, as stored in PTB-XL. We call this input raw: the digital signal as distributed in PTB-XL (after the
recording devices' own processing and the dataset authors' resampling), with no filtering, baseline correction, beat
detection or feature extraction by us. Three decimals in millivolts keep PTB-XL's resolution of 1\,\textmu V. The file was uploaded through the API's Files interface and attached to the request,
which made it available in the sandbox. The user message named the file and gave the patient's age and sex (ages of 90
years or more, which PTB-XL codes as 300, as ``90 or older''); Supplementary Section~\ref{app:example} shows one
message. No derived measurement, filter, report or label was provided.

\subsection{Model, sandbox, prompt and inference}\label{sec:inference}
We used \ModelID{} through the Anthropic Messages API with the server-side code-execution tool
(\CodeTool): a sandboxed container in which the model can run shell commands and Python (with NumPy, SciPy and pandas)
on the attached file, without internet access~\citep{anthropic2026codeexec}. Within one request, each code run is one round: the model
writes code, the sandbox executes it and returns the printed output, and the model then either runs more code or
answers. The prompt asked for measurement before judgement but did not set the number of runs. The API pauses a
request after a set number of tool rounds, and the protocol allowed up to two continuations in the same container. In the test run, each
request ran in a container of its own. Reasoning effort was \Effort{}, the output limit was \MaxTokens{} tokens per
API call (including reasoning, code and answer), and adaptive thinking was used with summarized display, so the API
returned a summary of the model's reasoning; in the test run this summary was empty for many requests, although every
request had used reasoning tokens (see (viii) under ``Post hoc analyses''). The model exposes no sampling parameters
(no temperature or seed), so outputs are not deterministic. The provider controls the sandbox image; error tracebacks
show Python 3.11, and package versions were not recorded. The model's published training-data cutoff is June
2026~\citep{anthropic2026models}. Development runs were the trial runs on validation-fold records in which prompts were
compared and the frozen prompt was chosen (``Development and freeze''); the test run was the single run of the frozen
prompt on the test fold. Development runs were made from \DevFirstSubmitUTC{} UTC on
\DevFirstSubmitDateLong{} to \DevLastCollectUTC{} UTC on \DevLastCollectDateLong, and the test run between
\TestFirstSubmitUTC{} and \TestLastCollectUTC{} UTC on \TestRunDateLong.

The system prompt (Supplementary Section~\ref{app:prompt}, SHA-256 prefix \PromptSHA) assigns the role of a
cardiologist with a Python sandbox, describes the file, and asks the model to work in two steps. In the first, it
filters the signal, detects beats, forms a median beat per lead with the baseline set from the PR segment, and
computes and prints a complete measurement table (rhythm, rate and regularity; PR, QRS duration from the earliest
onset to the latest offset over all leads, QT, and QTc by Bazett's formula; frontal axis; per lead, Q-wave duration
and depth, R, S and R$'$ amplitudes, ST level at the J point and 60\,ms later and T amplitude; the
Sokolow--Lyon~\citep{Sokolow1949} and Cornell~\citep{Casale1985,Casale1987} voltages and R in aVL; P-wave amplitude,
duration and terminal negativity in V1), and re-measures implausible values. In the second, it applies criteria for
each superclass to the table, weighing the overall pattern. The criteria were written for this study as a condensed
list of commonly used thresholds, not taken from a single source; Supplementary Table~\ref{tab:criteria} gives the
closest published criterion for each, from the AHA/ACCF/HRS recommendations for the interpretation of the
ECG~\citep{surawicz2009recommendations,rautaharju2009recommendations,hancock2009recommendations,wagner2009recommendations},
the fourth universal definition of myocardial infarction~\citep{thygesen2018fourth}, the guideline on bradycardia and
conduction delay~\citep{kusumoto2019bradycardia} and the original voltage criteria~\citep{Sokolow1949,Casale1985},
and marks where the prompt departs from it. The prompt defines the five superclasses as
in PTB-XL, states that rhythm findings alone are not superclass labels, and says that discrimination is measured, so
that probabilities should reflect graded evidence. It contains no class prevalences and no examples. The output was
constrained by a JSON schema (Supplementary Section~\ref{app:schema}): the record identifier, the key findings (free
text, which the prompt limits to at most two sentences naming the decisive findings and measurements) and five
independent probabilities.

Each record was sent as its own request through the Message Batches interface. The system prompt was identical across
requests and marked for prompt caching. The API reported the token counts of every request, from which cost was
computed at the provider's list prices (Supplementary Section~\ref{app:cost}).

\paragraph{Blinding.} Records were shuffled (seed 7) and given anonymous identifiers, which were also the file names;
the mapping to PTB-XL identifiers was stored locally and was never part of any request in the development or test
runs. The model saw no labels, reports or PTB-XL identifiers of the records it classified.

\subsection{Development and freeze}
The configuration was developed on the validation fold (fold 9) in \NDevRuns{} development runs. Each prompt was run
once on each set of records, and prompts were compared by macro AUROC against the validation labels:
\begin{enumerate}
  \setlength{\itemsep}{0pt}\setlength{\parsep}{0pt}
  \item A first, basic prompt asked the model to analyse the ECG as a cardiologist would and to measure whatever it
  judged necessary, without a measurement list or diagnostic criteria. It was run on 500 fold-9 records drawn at
  random (\texttt{protocol/dev500\_ids.txt}).
  \item The frozen prompt replaced this with the two steps described above: a fixed list of measurements to compute
  and print, then written criteria for each superclass. On the same 500 records it reached \DevVTwo{} (95\% CI
  \DevVTwoLow--\DevVTwoHigh). This was the only change that gave a measurable gain.
  \item Three variants of the frozen prompt were each compared with it on the same records by a paired bootstrap
  (Table~\ref{tab:development}): (a) ten labelled example recordings attached as files with their labels, two per
  superclass, taken from the training folds (1--8) so that no validation or test record served as an example (on 100
  of the 500 records, drawn at random); (b) probabilities for PTB-XL's 44 diagnostic statements instead of the five
  superclasses, aggregated to superclasses by their maximum (the same 100 records); and (c), after an error analysis
  on the 500 records, added rules for PTB-XL's labelling conventions (on 200 fold-9 records not used before, on which
  the frozen prompt reached \DevFreshVTwo, 95\% CI \DevFreshVTwoLow--\DevFreshVTwoHigh). None of them improved on the
  frozen prompt.
\end{enumerate}
The frozen prompt was kept because of the gain in step 2, because no variant improved on it, and because it needs no
labelled examples; it contains no training-fold data. All prompts are in \texttt{prompts/}. Development used
\NDevRecordsUsed{} distinct fold-9 records (the 500-record sample, the 200 new records and \NPilotExtra{} further
records from a \NPilot-record pilot whose other \NPilotInMain{} were in the 500-record sample) and cost
USD \DevCostTotal.

\begin{table}[tbp]
  \centering
  \caption{\textbf{Development on the validation fold (fold 9).} Each variant of the frozen prompt was run on the same
  records as the frozen prompt; macro AUROC and paired differences (variant minus frozen prompt) with 95\% percentile
  bootstrap CIs (\NBoot{} resamples). The 100 records are a random subset of the 500-record sample; on all 500, the
  frozen prompt reached \DevVTwo{} (95\% CI \DevVTwoLow--\DevVTwoHigh). The 200 records had not been used in
  development before. Each prompt was run once on each set of records; the CIs come from resampling the records
  (bootstrap), not from repeated runs. All runs were made before the freeze. The difference in the first row (the frozen prompt with the example files)
  was computed after the test run (post hoc).}
  \label{tab:development}
  \small
  \begin{tabular}{lrccc}
    \toprule
     & & \multicolumn{2}{c}{Macro AUROC} & \\
    \cmidrule(lr){3-4}
    Change & Records & Frozen prompt & Variant & Difference [95\% CI] \\
    \midrule
    + 10 labelled example files & 100 & \DevVTwoHundred & \DevBoth &
      $\DevBothDiff$ [$\DevBothDiffLow$, $\DevBothDiffHigh$] \\
    Statement-level output & 100 & \DevVTwoHundred & \DevStatements &
      $\DevStatementsDiff$ [$\DevStatementsDiffLow$, $\DevStatementsDiffHigh$] \\
    + rules for labelling conventions & 200 new & \DevFreshVTwo & \DevFreshVFive &
      $\DevFreshVFiveDiff$ [$\DevFreshVFiveDiffLow$, $\DevFreshVFiveDiffHigh$] \\
    \bottomrule
  \end{tabular}
\end{table}

The code, configuration and prompts were then frozen together with a written protocol that records
the prompt's SHA-256 hash, the endpoints and the analysis plan. The software refuses to build a test-fold run unless
the code, configuration and prompts are unchanged since the freeze and no experimental setting other than the
record-shuffle seed differs from the frozen configuration; the test-run manifest records no difference from the frozen
configuration, including the seed (7). The protocol was not registered externally.

\subsection{Test run and failure handling}
The \NTest{} test requests were sent in \NBatches{} batches of at most 250, one after another, with pre-specified stop
rules (a batch costing more than USD 20 or with more than 10\% failed answers). An answer had to contain five finite
probabilities; failed requests were to be retried once, requests paused by the API resumed up to twice, and records
still without an answer were to be given the training-fold class prevalences, with their number reported. None of
these was needed: no stop rule was triggered, and every request in the \NBatches{} batches returned a valid answer at
the first attempt. Every paid request was logged with its token counts and cost.

\subsection{Endpoints and statistical analysis}\label{sec:endpoints}
The pre-specified primary endpoint was macro AUROC over the five superclasses (scikit-learn \texttt{roc\_auc\_score},
macro average), the benchmark's metric~\citep{Strodthoff2021}, on all \NTest{} test records from the single test
run, with a 95\% percentile bootstrap confidence interval (CI)~\citep{Efron1993} (\NBoot{} resamples of records, seed 0,
keeping only resamples with positives and negatives of every class). Pre-specified secondary endpoints were per-class
AUROC, macro average precision (AUPRC), code runs, tokens and cost per record, and paired differences in macro AUROC,
computed on the same bootstrap resamples, against a logistic regression per class on age, age squared and sex (ages
coded 300 entered as 90), fitted on folds 1--8. Results of published models were compared as reported
by their authors, without significance tests. Per-class confidence intervals and per-class differences from the control use the same resamples and were not
pre-specified.

\subsection{Post hoc analyses}\label{sec:posthoc}
After the test run we added the following analyses.
\begin{enumerate}
  \renewcommand{\labelenumi}{(\roman{enumi})}
  \setlength{\itemsep}{0pt}\setlength{\parsep}{0pt}
  \item What the model did in the sandbox, from the transcripts: the non-empty lines of code per record; the code runs
  that ended with an error (a non-zero exit status, or an uncaught Python exception other than a broken pipe in the
  error output, since piped commands return the exit status of the last command) and, in records with an error,
  whether the last code run succeeded; whether the code called a filter (and its cutoffs), a peak detector and a
  median function; and whether the printed output contained the labels of heart rate, QRS, PR, QT or QTc, axis and
  Sokolow--Lyon voltage, a per-lead table of R and S amplitudes, or lines of 20 or more numbers. A printed label shows
  that a labelled value was printed, not where a stated value came from.
  \item Discrimination by diagnostic statement, for statements with at least \PHMinStatement{} test records, as the
  AUROC of the superclass probability for records with the statement against records without the superclass, with
  95\% CIs from \PHNStmtBoot{} stratified bootstrap resamples; this measures how well the superclass probability
  separated these records, not whether the statement itself was identified.
  \item Subgroups by sex, age (under 50 years, 50--69 years, 70 years or older), rhythm (atrial fibrillation or
  flutter), signal-quality flags and extra beats noted in the PTB-XL metadata, and recording device (the device values
  as recorded in PTB-XL, with at least \PHMinSubgroup{} test records); for the rhythm subgroups also as the mean AUROC
  of the four classes other than NORM, which had few positive records in atrial fibrillation or flutter.
  \item Calibration, that is, whether the probabilities can be read at face value (of records given a probability of
  about 0.3 for a class, about 30\% should have it). AUROC depends only on how records are ranked, so it does not
  measure this. We report the expected calibration error (ECE: the absolute gap between mean predicted probability and
  observed frequency in ten equal-width probability bins, averaged with weights by the number of records per bin), the
  Brier score (the mean squared difference between probability and label) and the calibration slope (the slope of a
  logistic regression of the label on the log-odds of the probability: 1 for calibrated probabilities, above 1 for
  probabilities that are not extreme enough).
  \item MI discrimination by the likelihood of the MI statements: each MI-positive record was grouped by the highest
  likelihood among its MI statements (100; 50--80; 15--35) and compared with all test records without MI, with 95\% CIs
  from \PHNStmtBoot{} stratified bootstrap resamples as in (ii); the age-and-sex control was evaluated on the same
  groups.
  \item A bootstrap that resampled patients with replacement instead of records, each drawn patient contributing all
  of their test records (\NBoot{} resamples), because records of the same patient are not independent.
  \item The measurements stated in the key findings, that is, the values that the model wrote in the answer's free
  text, whatever their source: heart rate, PR, QRS duration, QTc, axis and Sokolow--Lyon voltage. The key findings of
  the record in Fig.~\ref{fig:example}e, for example, state all six, among them an axis near $+\ExampleAxisStated^\circ$,
  whereas the code had printed an axis of $\ExampleAxisPrinted^\circ$ (Fig.~\ref{fig:example}c). They were extracted
  from the key-findings text with regular expressions (the first value with its unit after the
  measurement's label within its clause; the midpoint of a stated range; for Sokolow--Lyon the stated sum, not a
  component; details in Supplementary Table~\ref{tab:measurements}), and mentions of atrial fibrillation or flutter
  were detected with a pattern that excluded negations. These were tabulated against the PTB-XL statements and the
  heart-axis field of the PTB-XL metadata, with the criteria of the frozen prompt where it gives one (QRS 120\,ms or
  more for complete bundle branch block and non-specific intraventricular conduction disturbance, 110--119\,ms for
  incomplete right bundle branch block, PR over 200\,ms for first-degree atrioventricular block, Sokolow--Lyon voltage
  3.5\,mV or more for left ventricular hypertrophy) and conventional limits otherwise. The prompt gives no angle for
  fascicular block, so the stated axis in left anterior fascicular block was compared with two thresholds, $-30^\circ$
  (left axis deviation) and $-45^\circ$ (the AHA/ACCF/HRS criterion~\citep{surawicz2009recommendations}).
  \item Mentions of PTB-XL by name in the reasoning summaries, among the \PHSummariesNonEmptyN{} of \NTest{} test
  records whose summary was not empty.
  \item Two analyses on the validation fold: the age-and-sex control on the \NDevVTwoAll{} records read with the frozen
  prompt and on the whole fold, and the paired difference of the example-file variant from the frozen prompt (first
  row of Table~\ref{tab:development}), computed from development runs made before the freeze.
  \item A local replay of the model's code: every code run of the test run was executed again, record by record and
  in the original order, on a local computer without network access, in a sandbox (bubblewrap) that mirrored the API
  container: the record's file was in the directory named by the same environment variable, and working files were
  kept across the runs of one record. It used Python \RPPython{} with NumPy \RPNumPy, SciPy \RPSciPy, pandas
  \RPPandas{} and matplotlib \RPMatplotlib, chosen to match what the tracebacks and printed output show (Python 3.11,
  NumPy 2). The printed output of each run was compared with the output the API returned and classed as identical,
  equal up to rounding (the same text, with every number within one unit of its last printed digit or within a
  relative $10^{-9}$) or different; an error as defined in (i) counted as reproduced if the local run also ended with
  an error of the same exception type. The model was not queried again, so the replay tests the code runs, not the
  model's answers.
\end{enumerate}

\subsection{Software}
Local analyses used Python 3.13.9 with NumPy 2.3.5, SciPy 1.16.3, pandas 2.2.3, scikit-learn 1.7.2, wfdb 4.3.1,
matplotlib 3.10.6 and the Anthropic Python SDK 0.117.0; the exact versions of the packages that the code uses directly
are pinned in \texttt{requirements-lock.txt}.

\subsection{Ethics}
This is a secondary analysis of PTB-XL~\citep{Wagner2020,Wagner2022physionet}, a public, de-identified dataset; no new
data were collected and no participant was contacted. The PTB-XL authors report that their institutional ethics
committee approved the publication of the anonymous data in an open-access database (PTB-2020-1), and that personal
information in the metadata was pseudonymized and recording dates were shifted by a random offset per
patient~\citep{Wagner2020,Wagner2022physionet}. This analysis therefore required no further ethics approval, and none
was sought. The de-identified recordings were processed through a commercial API (Anthropic).
